% THIS HAS BEEN SAID ELSEWHERE AND IN CONCLUSIONS \textbf{We should say that the proposed rule-based J-LLM benchmark is a much better benchmark than even compile based benchmark compile@1, also much better than similarity based metrics like CodeBLEU and ROUGE-LSUM. The J-LLM benchmark is much closer to the ultimate performance metric pass@1, which is the real performance of the S-LLMs. The J-LLM benchmark is much more efficient than execution based benchmarks, as simulation tasks are expensive to run. A complex simulation task like high-fidelity computational fluid dynamics (CFD) can take months to run. The J-LLM benchmark is a much better way to evaluate the performance of S-LLMs in generating DTs. The J-LLM benchmark uses Reference code and API documentation in the determination process, which may increase the API cost, but the cost of the API is acceptable. We should grow the dataset used to improve our J-LLM, including more dynamic systems, e.g., fluid-solid interaction; multi-phase flow. The rule-based J-LLM can be a better judge compared to other metrics, as it is more closed to the pass@1, which shows the real and ultimate performance of the S-LLMs for simulation tasks.}
% THIS HAS BEEN SAID ELSEWHERE AND IN CONCLUSIONS \textbf{Image, you can know the code by reading even better than by runing the code itself.
%We also should say that we build a really good quality and difficult simulation/DT dataset, although most of LLMs are pretrain on, even the most advanced S-LLMs only have less than $13\%$ pass@1.  }

\PARstart{I}{t} is only a matter of time before a non-expert will be able to interact with an S-LLM in order to produce a sophisticated digital twin (DT) for a physical system of interest. The DT request might read ``generate a model of a VIPER rover operating on regolith and equipped with two monocular cameras, one stereo camera, and an autonomy stack that uses sensor feeds to ensure obstacle avoidance while operating in lunar conditions.'' Such a model can be generated for a target simulator should an expert dedicate enough time to the task. However, the interest here lies in the scenario where the DT generation falls upon an LLM. Specifically, this contribution proposes a benchmark that measures the proficiency of such an S-LLM when it comes to generating DTs for a target simulator. 

We work under the assumption that a DT assumes the form of a simulator-ready script or a collection of configuration files that enable a virtual experiment via computer simulation. In the VIPER rover example above, the DT is instantiated through hundreds of lines of Python/C++ code or a proprietary modeling language. Alternatively, the DT may be specified via structured configuration files (e.g., JSON), which can likewise span hundreds to thousands of lines. For instance, the rover definition must specify, for each wheel, the center of mass, mass moments of inertia, initial pose and velocity, collision and visualization meshes, and frictional parameters; analogous properties must be defined for the chassis. Camera models require parameters spanning optics and photometry (e.g., aperture, exposure, focal length, BRDF-related quantities, and image-signal-processing settings). The terrain model may involve bulk density, internal friction angle, cohesion, and particle-size distribution. Finally, an autonomy stack may be integrated by configuring a ROS2 environment and selecting perception, planning, and control components. Ultimately, the objective is for an S-LLM to generate the DT construction artifact (e.g., Python code or configuration files); when the user request is underspecified, the S-LLM should either adopt principled defaults or elicit the missing information through targeted clarification.


Analyzing complex multi-physics phenomena via simulation requires sophisticated DTs. 
S-LLMs can potentially handle these through multiple rounds of interaction, 
incorporating natural language feedback from humans. 
To the best of our knowledge, dedicated evaluation protocols for the task of generating DTs are nonexistent, 
as generic similarity-based metrics like CodeBLEU \cite{ren2020codebleu} and ROUGE-L \cite{lin-2004-rouge} 
do not align well with real-world simulation applications and can not capture the complexity of the simulation tasks.

Furthermore, execution-based benchmarks such as Pass@k \cite{chen2021evaluating} can be overly stringent in simulation-centric settings: a single minor defect (e.g., an omitted parameter, setting up an inadequate time step convergence criteria, or swapping value of the time step for that of the relative solution tolerance) may cause a total test failure, yielding a zero score despite the generated digital twin being qualitatively close to correct. More importantly, designing high-coverage and reliable unit tests for complex physics simulations is inherently difficult and rarely scalable. Unlike LeetCode-style problems, where inputs, outputs, and edge cases can be precisely specified, simulation tasks involve long-horizon dynamics, continuous states, numerical sensitivity, and complex scenario logic, making exhaustive, stable, and reusable unit-test suites impractical to author and maintain. As a result, execution-based evaluation becomes costly to extend across diverse systems, scenarios, and multi-turn iterations. It also makes downstream approaches that depend on test-verifiable reward signals (e.g., reinforcement learning with verifiable rewards) substantially less tractable, since they would require scenario-specific test design to provide reliable, dense supervision.

Most importantly, conventional execution-based metrics (e.g., Compile@k and Pass@k)
typically collapse evaluation into a single scalar, 
offering little diagnostic insight into what is wrong and why it is wrong in the DT generated by the S-LLM.
This lack of interpretability is particularly problematic for simulation-grade DTs,
which often comprise hundreds to thousands of lines of code. Indeed, a single scalar score provides limited guidance for debugging, model improvement, or iterative refinement.
In practice, what is needed is an evaluation signal that is both reliable and interpretable, i.e., one that can
attribute errors to concrete aspects of the generated artifact (e.g., missing components,
incorrect parameterization, inconsistent physics assumptions, or incomplete scenario logic),
and provide actionable feedback.

Against this backdrop, the goal of this contribution is to establish an interpretable benchmark for assessing how well
simulation-focused LLMs (S-LLMs) generate DTs. Beyond serving as an evaluation metric,
such an interpretable scoring signal can also support future learning paradigms that require dense and
informative rewards, e.g., improving an S-LLM via training objectives that do not rely on writing
exhaustive unit-test suites for every scenario.
        
Insofar as the nomenclature used in the contribution is concerned, an analysis task in the real world corresponds to a simulation scenario in the digital world. A physical system, e.g., the VIPER rover, has a digital twin. Several digital twins, e.g., those of the rover, terrain, sensor suite, are combined in a simulation scenario towards running a virtual experiment.
 In the absence of a sim-to-real gap, the results collected in simulation are representative of results that
  would have been collected had the experiment taken place in the real world.
   Here, the expectation is that an S-LLM is used to generate the required digital twin[s],
    place them in a simulation scenario, and run a virtual test of interest in a simulator of choice.
     This contribution is not concerned with how a good S-LLM should be set up.
      Rather, the interest is in measuring the ability of an S-LLM to generate good DTs;
       i.e., in benchmarking the S-LLM's performance. As such, the primary interest is in the J-LLM, 
and without loss of generality for the methodology discussed, the simulator used to introduce SimBench is
 PyChrono \cite{pychrono2022}, the Python wrapper of the multi-physics Chrono simulator \cite{chronoOverview2016}. 

The PyChrono choice is justified by the authors' familiarity with the simulator, its broad use in the CAE community, and the fact that it is open source and available on GitHub \cite{projectChronoGithub,projectChronoWebSite}. PyChrono, which was released in 2022 and provides Python bindings to the core C++ Chrono library, has been downloaded from Conda-forge more than 160,000 times \cite{chronodownload}. The Chrono User Forum has more than 800 users \cite{projectChronoForum}. Chrono::Engine provides core functionality for multibody dynamics and nonlinear finite element analysis in the presence of friction and contact. Chrono::Vehicle provides a collection of wheeled and tracked vehicles, facilitating high-fidelity vehicle dynamics analysis, engine and driveline simulation, and deformable terrain modeling. Chrono supports fluid-solid interaction and granular dynamics simulation. Chrono::ROS provides a ROS bridge that enables Chrono to interact with third party autonomy stacks \cite{ROS-2009}. Chrono::Sensor provides sensor modeling and simulations support. Several other modules provide ancillary support for interfacing to SolidWorks, MATLAB, and other third party solvers via FMU/FMI interfacing \cite{FMI_3.0}.

This work makes the following contributions:

(i) \textbf{SimBench: a multi-turn, interpretable benchmark for DT generation.} We introduce SimBench, a benchmark methodology designed to assess an LLM's ability to produce \emph{simulator-ready} digital twins (DTs) within a real simulation ecosystem via multi-turn code editing. Relative to prior code-generation and simulation benchmarks, SimBench explicitly combines (1) multi-turn interaction, (2) rubric-grounded, interpretable scoring with reference/document conditioning, and (3) partial execution anchoring to connect model outputs to runtime behavior. As DT technology becomes increasingly relevant across engineering and scientific domains \cite{digitalTwinReportNAcademies2023}, automating DT construction is highly desirable; a principled, transparent, and interpretable evaluation protocol represents a step toward developing effective simulation-focused LLMs (S-LLMs).

(ii) \textbf{A curated expert dataset and a calibrated J-LLM evaluation pipeline.} We develop a curated dataset authored by simulation experts, consisting of 102 exemplar demonstrations spanning 34 distinct physical systems. For each system, we provide three demonstrations with progressively increasing complexity, together with instructions that guide the transformation from simpler to more intricate configurations. This dataset serves both as (a) a benchmark substrate for standardized evaluation, and (b) a template for scaling data generation for supervised fine-tuning. Using SimBench, we evaluate 33 S-LLMs and collect over 3{,}000 multi-turn dialogues, each paired with expert-assessed scores. This collection can further serve as a preference-style dataset for training or calibrating judge models (J-LLMs) and/or reward models that provide interpretable supervision without requiring scenario-specific unit-test engineering. All datasets, benchmarks, and evaluation code are open-sourced.\footnote{\url{https://github.com/uwsbel/SimBench}}

(iii) \textbf{Generality beyond Chrono.} While SimBench is instantiated for Chrono-based DT generation, the methodology is simulator-agnostic. Given an analogous expert dataset for another simulation platform (e.g., OpenFOAM \cite{openFoam} or PyBullet \cite{bulletPhysicsEngine2020}), the same workflow can be used to construct a calibrated J-LLM suitable for judging DT artifacts in that target ecosystem.


\begin{figure*}[htbp]
    \centering
    \includegraphics[width=0.875\textwidth]{pipeline.pdf}
    \caption{The SimBench pipeline for evaluating S-LLMs. The J-LLM is calibrated using a validation set containing pairs of ground truth and generated DTs. The prompts given to the J-LLM are iteratively optimized to match the score provided by the expert. Then the J-LLM is used to evaluate the S-LLM based on the generated DTs, ground truth DTs, and the API documentation.}
    \label{fig:your-label}
\end{figure*}
